\chapter{Cadre général du projet}


% Une section

% Exemple d'une section qui porte une référence à une bibliographie
% NB: il faut bien suivre le syntaxe pour ne pas tomber dans le cas où il y a une référence dans la table des matières.

\section*{Introduction}

Le premier chapitre présente une étude préliminaire du projet. Nous commencerons par présenter l'organisme d'accueil, en mettant en évidence son rôle et son contexte dans le domaine d'étude. Ensuite, nous examinerons la problématique à laquelle notre projet répond, ainsi que les objectifs clés que nous visons. Enfin, nous décrirons la méthodologie de travail, en mettant l'accent sur les approches agiles telles que Scrum et la méthodologie de conception que nous avons utilisée.

\section{Présentation de l’organisme d’accueil } 
Nous introduisons dans cette partie l’historique d’ASTEELFLASH , ses activités et ses différents directions.
\subsection{Présentation d'ASTEELFLASH TUNISIE  }
Asteelflash Tunisie est une filiale du groupe français ASTEELFLASH. Elle a été fondée en 2010 à Megrine par M. Gilles Benhamou, fondateur du groupe français ASTEELFLASH. En 2016, le site tunisien a été transféré dans une nouvelle usine située à La Soukra, à Tunis. Avec une superficie de 15 000 m², cette nouvelle usine est le cœur du projet Excellence. Elle a rangé à positionner ASTEELFLASH parmi les meilleurs fournisseurs mondiaux de services de fabrication électronique (EMS) en mettant en place un outil industriel solide et moderne, doté d'une technologie de pointe, soutenu par des certifications de haut niveau, une politique de formation et des ressources conséquentes.

Le site de production de l'usine de La Soukra est un endroit idéal pour la fabrication et l'assemblage de projets à bas coût. Il dessert les pionniers de l'industrie européenne dans ces segments, en offrant des délais de production courts par rapport à d'autres régions du monde. L'établissement de La Soukra a développé un savoir-faire spécifique pour répondre aux exigences les plus élevées des industries, avec des équipements de pointe conformes aux normes industrielles les plus strictes.

\newpage

\begin{longtable}{|m{4cm}|m{8cm}|}
                
                \hline 
                    \textbf{Raison sociale} & \textbf{ASTEELFLASH TUNISIE}\\
                \hline
                
                    Date de création
                    & 2000 \\
            
            
                \hline
                    Forme juridique
                    & Société À Responsabilité Limitée (SARL) \\
                    
                    
                \hline
                    Siège social
                    & 104 Avenue de l’UMA La Soukra Ariana Tunisie \\
                               
                \hline
                    Tél
                    & +216 36 066 015 \\
                               
                \hline
                    Fax
                    & +216 36 066 000 \\

                \hline
                    Site web
                    & www.asteelflash.com \\

                \hline
                    Superficie
                    & 15.000 m² d'espace de production \\

                \hline
                    Propriétaire
                    & Gilles Benhamou \\

                \hline
                    Effectif
                    & 1000 employés \\

                \hline
                    Régime de travail
                    & 40 heures /semaine (3 postes) \\

                \hline
                    Certification
                    & -ISO 9001 \newline
                      -IATF 16949 (ex. ISO/TS 16949) \newline
                      -ISO 14001\\
                               

                \hline

              \captionsetup{justification=centering,margin=2cm}
              \caption{Fiche technique de l'entreprise}

            \end{longtable}
\subsection{Organigramme de l'ASTEELFLASH}
La figure 1.1 représente la hiérarchie de l’ASTEELFLASH avec ses différents
départements.

    \begin{figure}[h!]
    \centering
    \includegraphics[scale=0.7]{Organigrame.PNG}\\[0.5cm]
     \caption{Hiérarchie de Département IT d'ASTEELFLASH }
    \end{figure}\
    
    \subsection{Présentation du groupe ASTEELFLASH }
 Le groupe Asteelflash est un sous-traitant électronique français crée en 1999 par M Gilles Benhamou, à Neuilly Plaisance, à proximité de Paris. 
 Classé parmi les deux premières EMS (Electronic Manufacturing Services) européennes et les quinze premières mondiales, le groupe compte 23 sites de production sur 4 continents avec un effectif de plus de 5 700 collaborateurs. 
 
Dès sa création l’entreprise a fixé des objectifs ambitieux et se lança dès ses premiers pas dans une stratégie fortement axée sur l’initiative de développement tout en favorisant des partenariats solides avec ses clients. Elle utilise sa créativité et son savoir-faire dans la fabrication et l’intégration des produits très complexes, allant de la conception initiale du produit jusqu’à la production en série, y compris le prototypage, le traitement des commandes directes et les services après-vente. La chaîne innovante d'approvisionnement, la technologie et la qualité des solutions proposées permettent à Asteelflash de réduire, sur une base continue, les coûts et les délais de commercialisation au bénéfice de tous ses clients. 

Asteelflash sert des entreprises de haute technologie dans le monde entier appartenant à un large éventail de segments de marché, incluant les secteurs industriels, la gestion énergétique, le traitement des données, la défense et l'aérospatiale, le transport (automobile et ferroviaire) ainsi que le secteur médical.
        
\subsection{Activités et produits}
    \begin{itemize}
        
    \item [$-$]\textbf {Pour le pôle électronique, ASTEELFLASH a comme activités :}

    \begin{itemize}[label=\textbullet]
    \item Câblages des cart
    \item Pose composants CMS
    \item Double refusions
    \item Programmation du test
    \item Assemblages de sous-ensemble électroniques
    \end{itemize}
   
    \item [$-$]\textbf { Pour le pôle plasturgie (électromécanique) :}
   
    \begin{itemize}[label=\textbullet]
    \item Surmoulage et injection
    \item Assemblage électromécanique et intégration
    \end{itemize}
    
    \item [$-$]\textbf { ASTEELFLASH produit plusieurs cartes électroniques de différents clients entre autres on cite:}
  \begin{itemize}[label=\textbullet]
    \item Cartes compteur d’énergie (SOCOMEC)
    \item Cartes thermostats (GROUPE MULLER, NOIROT)
    \item Cartes de commande (FAGOR-BRANDT)
    \item Cartes minuteries (HAGER)
    \item Cartes de régulation (DE DIETRICH)
    \end{itemize}
      \end{itemize}
\subsection{Les différentes directions}
   
    Les principales directions d’ASTEELFLASH sont : 
   
    \begin{itemize}
      
    \item [$-$]\textbf{Direction Générale :}
    
    \begin{itemize}[label=\textbullet]
    \item Piloter l'activité et le développement de l'entreprise.
    \item Coordonner les missions de conseil en organisation et management. 
    \end{itemize}

     \item [$-$]\textbf{Direction Qualité :}
   
    \begin{itemize}[label=\textbullet,]
    \item Définir et mettre en œuvre la politique qualité de l'entreprise.
    \item Coordonner les activités de pilotage et de surveillance de la performance des procédures qualité. 
    \end{itemize}

  
    \item [$-$] \textbf{Direction Engineering:}    
    \begin{itemize}[label=\textbullet,font=\normalsize]
    \item Diriger l’ensemble des équipes produisant les études techniques.
    \item Anticiper les impacts des transformations et conduire l’accompagnement au changement. 
    \end{itemize}

    
     \item [$-$]\textbf{Direction Ressources Humaines:}
    
    \begin{itemize}[label=\textbullet]
    \item Définir, valider, piloter et suivre la réalisation de la stratégie RH de l’entreprise.
    \item Assurer le lien entre la direction générale et le personnel avec prise en charge des relations sociales au sein de l'organisation.  
    \end{itemize}

     \item [$-$]\textbf{Direction Achats:}

    \begin{itemize}[label=\textbullet]
    \item Définir la stratégie achat en collaboration avec la direction générale de l'entreprise et veille à sa mise en œuvre à travers un choix ciblé de produits/services et fournisseurs.
    \item Négocier les contrats d'achats et coordonner la chaîne logistique.   
    \end{itemize}

   
     \item [$-$]\textbf{Direction Production:}
    
    \begin{itemize}[label=\textbullet]
    \item Superviser la fabrication de gammes de produits.
    \item Veiller à l'utilisation optimale des ressources.  
    \item Garantir le respect des contraintes (qualité, coûts, délais). 
    \item Superviser les cadences de production et piloter les actions d'amélioration continue.   
    \end{itemize}

  
    \item [$-$] \textbf{Direction Logistique :}
  
    \begin{itemize}[label=\textbullet,font=\normalsize]
    \item Coordonner et optimiser la circulation des matières premières et des produits.
    \item Élaborer la politique de gestion des flux et le planning de production.   
    \item Traiter l’approvisionnement des matières premières.  
    \item Optimiser le stockage et la manutention. 
    \item Assurer l’expédition des produits finis. 
    \end{itemize}

   
     \item [$-$]\textbf{Direction Lean :}
    
    \begin{itemize}[label=\textbullet,font=\normalsize]
    \item Optimiser les performances industrielles des projets d’amélioration et d’innovation.
    \item Réaliser des audits d’évaluation des performances industrielles et concevoir le plan d’action pour l’amélioration des processus industriels.
    \item Assurer le déploiement des solutions d’amélioration et piloter la conduite du changement.  
    \item Assurer la formation du personnel technique aux nouvelles méthodes. 
    \end{itemize}

  
    \item [$-$] \textbf{Direction Financière :}
    
    \begin{itemize}[label=\textbullet,font=\normalsize]
    \item Définir en collaboration avec la direction générale, la stratégie financière de l'entreprise.
    \item Planifier, gérer et coordonner les activités financières et comptables de l’entreprise.
    \end{itemize}
    
    \end{itemize}
    
\section{Etude de l’existant}

Le système de gestion du parc informatique qui repose sur des décharges papier et des fichiers Excel partagés, d’où il y aura toujours un risque de perdre ces décharges ou de modifier le fichiers Excel par les non connaisseurs. 
Actuellement, ces méthodes traditionnelles présentent plusieurs limitations et inefficacités. La gestion des actifs matériels et logiciels se fait de manière dispersée, ce qui rend difficile la surveillance et la maintenance adéquate du parc informatique. 
De plus, le suivi des licences, la planification des maintenances et la génération de rapports se font manuellement, ce qui entraîne des risques d'erreurs et une perte de temps considérable.

    
Parmi les vulnérabilités de l’utilisation des fichiers Excel :

\begin{itemize}[label=\textbullet,font=\normalsize]
    \item Difficulté de gestion de stock actuel
    \item Risque d’attribuer le même matériel à deux personnes différentes
    \item Contrainte de suivi des matériels affectés
    \item Utilisation gênante des papiers de décharge lors de l’affectation
    \item Pas de traçabilité pour la réception des matériels, nom du fournisseur, la quantité la qualité et la marque 
    \item Saisie tardive des entrées/sorties (n’est pas à temps)
\end{itemize}


\subsection{Diagnostic détaillé du processus actuel}

Afin de mieux comprendre les insuffisances du système existant, nous avons analysé le déroulement réel des principales opérations effectuées par l'équipe IT. Cette analyse a montré que les tâches de suivi sont réparties entre plusieurs supports : fichiers Excel, formulaires papier, échanges par courrier électronique et communications directes entre les techniciens. Cette dispersion rend la consolidation des informations difficile, surtout lorsque plusieurs personnes interviennent sur le même équipement.

Le processus actuel peut être résumé comme suit :

\begin{itemize}[label=\textbullet,font=\normalsize]
    \item Lorsqu'un équipement est reçu, ses informations sont saisies manuellement dans un fichier Excel. Cette saisie dépend de la disponibilité du technicien responsable et peut donc être retardée.
    \item Lorsqu'un équipement est affecté à un employé, une décharge papier est préparée puis archivée. Le lien entre la décharge et l'enregistrement Excel n'est pas toujours immédiat.
    \item Lorsqu'un équipement est déplacé, retourné ou envoyé en mission temporaire, l'information doit être mise à jour manuellement. En cas d'oubli, l'état réel du parc ne correspond plus à l'état indiqué dans le fichier.
    \item Lorsqu'un stock devient critique, l'alerte dépend de la consultation périodique du fichier. Il n'existe pas de notification automatique permettant d'anticiper les ruptures.
    \item Lorsqu'une question d'audit apparaît, l'équipe doit rechercher dans plusieurs fichiers, messages et documents papier pour reconstituer l'historique de l'équipement.
\end{itemize}

Cette organisation entraîne une dépendance forte à la mémoire des personnes et aux habitudes de travail. Elle peut fonctionner lorsque le volume d'équipements est faible, mais elle devient fragile dans un contexte industriel où les mouvements de matériel sont fréquents et où la disponibilité du système d'information est critique.

\subsection{Impacts organisationnels et techniques}

Les limites du système existant ne concernent pas uniquement la saisie des données. Elles ont également des conséquences directes sur la qualité du service IT, la sécurité du patrimoine informatique et la capacité de l'entreprise à prendre des décisions rapides.

\begin{longtable}{|m{4cm}|m{5cm}|m{5cm}|}
\hline
\textbf{Constat} & \textbf{Impact observé} & \textbf{Conséquence pour ASTEELFLASH} \\
\hline
\endhead
\endfoot
\endlastfoot
\hline
Données dispersées entre Excel et papier & Difficulté à obtenir une vision globale du parc & Temps perdu lors de la recherche d'information et risque de décision basée sur des données incomplètes \\
\hline
Absence de contrôle automatique des doublons & Possibilité d'affecter le même équipement à plusieurs personnes & Perte de fiabilité de l'inventaire et conflits lors de la restitution du matériel \\
\hline
Historique incomplet des mouvements & Traçabilité limitée des affectations, missions et retours & Difficulté à identifier la responsabilité en cas de perte, panne ou retard de restitution \\
\hline
Suivi manuel du stock & Détection tardive des seuils critiques & Risque d'indisponibilité de matériel pour les nouveaux collaborateurs ou les besoins urgents \\
\hline
Absence de centralisation des fournisseurs et commandes & Recherche lente des informations d'achat et de garantie & Gestion moins efficace des achats, des réceptions et des renouvellements \\
\hline
Accès non contrôlé aux fichiers partagés & Modification accidentelle ou non autorisée des données & Risque d'altération des informations sensibles du parc informatique \\
\hline
\captionsetup{justification=centering,margin=1cm}
\caption{Synthèse des impacts du système existant}
\end{longtable}

À partir de cette analyse, nous constatons que le besoin principal ne se limite pas à numériser un fichier Excel. Il s'agit plutôt de mettre en place un véritable système de gestion du parc informatique, capable de centraliser les données, d'automatiser les contrôles, d'assurer la traçabilité et d'accompagner les décisions opérationnelles.
\begin{figure}[h!]
    \centering
    \includegraphics{affectation.PNG}\\
     \caption{Fichier Excel D’affectation des équipements}
\end{figure}

L’idée générale du projet consiste à concevoir un outil applicatif qui pourra de façon concrète aux équipes IT de ASTEELFLASH de gérer son parc. La gestion de parc informatique, c'est le suivi en temps réel du patrimoine informatique, matériel et logiciel de l'entreprise. Elle offre une vision globale de l'état, du suivi et des coûts des appareils utilisés dans l'entreprise. Il s'agit non seulement de recenser les différents types de machines présentes dans l'entreprise, leur nombre, leur localisation et les logiciels qui y sont installés mais de connaître aussi l'année d’acquisition de ces matériels et logiciels, le nom des fournisseurs etc.

\section{Cadre générale du projet}
Nous expliquons dans cette partie, la problématique et nous présentons la solution proposée.
\subsection{Problématique}
La gestion d'un parc informatique constitue souvent un défi majeur pour les organisations. Les ressources informatiques peuvent être dispersées dans différents départements, rendant ainsi leur suivi et leur maintenance difficiles. De plus, l'absence d'un système centralisé peut entraîner des problèmes tels que la perte de matériel, les licences expirées, les mises à jour manquantes et les performances médiocres des systèmes. Il est crucial de trouver une solution qui permette une gestion efficace du parc informatique, en fournissant un aperçu complet de tous les actifs, en automatisant les tâches récurrentes et en facilitant la prise de décisions éclairées.

Parmi les défis spécifiques rencontrés dans la gestion d'un parc informatique, il y a le manque de traçabilité et de suivi. De nos jours, le contrôle et le suivi des opérations techniques et financières dans les entreprises se font rarement de manière automatisée et sécurisée. Une problématique fréquente concerne les équipements non recensés. Dans la gestion d'un réseau d'entreprise, il peut être difficile de suivre les équipements en panne et de les remplacer par des équipements adéquats qui rempliraient exactement les mêmes fonctions (même système, même modèle, même périphérique). Ceci est dû au fait que le stock des ressources de l'entreprise n'est pas inventorié.

Un autre défi réside dans l'insécurité. Les utilisateurs ont tendance à utiliser les équipements de l'entreprise à des fins personnelles, ce qui peut entraîner la connexion de périphériques de stockage potentiellement infectés par des virus, l'installation de logiciels malveillants et la désactivation du pare-feu, entre autres. Une solution concrète doit être mise en place pour prendre en charge ces différents problèmes et assurer la sécurité des équipements de l'entreprise.

\subsection{Solution}
Pour contourner ces problèmes, nous proposons le développement d'une application web dédiée à la gestion du parc informatique. Cette solution offrira une interface conviviale et intuitive, permettant aux utilisateurs d'accéder facilement aux fonctionnalités clés. Grâce à cette application, les administrateurs informatiques pourront suivre en temps réel l'état du parc, gérer les actifs matériels et logiciels, planifier les tâches de maintenance, effectuer des inventaires et générer des rapports détaillés.

L'application permettra une gestion centralisée de tous les actifs informatiques, en fournissant un aperçu complet et précis de l'ensemble du parc. Les utilisateurs pourront visualiser les informations essentielles telles que les spécifications matérielles, les licences logicielles, les dates d'expiration, les mises à jour disponibles et les performances des systèmes. Ils pourront également localiser facilement les ressources dans l'organisation, ce qui facilitera la maintenance, la réparation ou le remplacement des équipements défectueux.

L'automatisation des tâches récurrentes sera un élément clé de cette application. Les utilisateurs pourront configurer des rappels pour les renouvellements de licences, les mises à jour logicielles et les opérations de maintenance préventive. Cela permettra de réduire les erreurs humaines et d'optimiser l'utilisation des ressources informatiques. De plus, l'application générera des rapports détaillés sur l'état du parc, les coûts de maintenance, les performances des systèmes, les risques potentiels et d'autres métriques pertinentes. Ces rapports faciliteront la prise de décisions éclairées en matière de planification, de budgétisation et d'optimisation du parc informatique.


\subsection{Objectifs détaillés de la solution}

La solution proposée vise à répondre à plusieurs objectifs complémentaires. Ces objectifs ont été définis à partir de la problématique, des contraintes de l'entreprise et des échanges avec les responsables du département IT.

\begin{itemize}[label=\textbullet,font=\normalsize]
    \item \textbf{Centraliser l'information :} regrouper les données relatives aux équipements, fournisseurs, commandes, affectations, missions et employés dans une base de données unique.
    \item \textbf{Fiabiliser les opérations :} réduire les erreurs de saisie grâce à des formulaires structurés, des listes de choix, des contrôles de validation et des messages d'erreur compréhensibles.
    \item \textbf{Assurer la traçabilité :} enregistrer automatiquement les opérations importantes afin de conserver un historique exploitable lors des audits ou des recherches.
    \item \textbf{Améliorer la réactivité :} permettre aux utilisateurs de rechercher rapidement un équipement, de connaître son état, son affectation et son historique.
    \item \textbf{Automatiser les alertes :} notifier les personnes concernées lorsque le stock atteint un seuil minimal ou lorsqu'une mission arrive à échéance.
    \item \textbf{Sécuriser l'accès :} intégrer l'authentification Active Directory afin d'utiliser les comptes existants de l'entreprise et de limiter l'accès aux membres autorisés.
    \item \textbf{Faciliter le pilotage :} fournir un tableau de bord permettant de visualiser la répartition du parc par type, état et disponibilité.
\end{itemize}

\subsection{Périmètre fonctionnel du projet}

Le périmètre du projet a été volontairement défini de manière progressive afin de livrer une solution exploitable dans le délai du PFE. Le premier périmètre couvre les fonctionnalités prioritaires nécessaires au département IT. Les évolutions plus avancées, comme la gestion complète des licences logicielles ou l'inventaire automatique réseau, sont identifiées comme perspectives.

\begin{longtable}{|m{4cm}|m{6cm}|m{4cm}|}
\hline
\textbf{Domaine} & \textbf{Fonctionnalités couvertes} & \textbf{Statut dans le projet} \\
\hline
\endhead
\endfoot
\endlastfoot
\hline
Authentification & Connexion LDAP, contrôle des groupes, distinction Administrateur/Utilisateur IT & Inclus \\
\hline
Gestion des équipements & Ajout, modification, consultation, suppression, recherche, suivi de l'état et du type & Inclus \\
\hline
Gestion du stock & Consultation par catégorie, seuil minimum, alerte automatique & Inclus \\
\hline
Affectations & Affectation d'un équipement à un employé, consultation de la liste et historique & Inclus \\
\hline
Missions temporaires & Création de mission, suivi des dates, rappel automatique par e-mail & Inclus \\
\hline
Fournisseurs et commandes & Gestion des fournisseurs, bons d'achat, réception des commandes & Inclus \\
\hline
Reporting & Tableau de bord statistique et indicateurs de disponibilité & Inclus \\
\hline
Licences logicielles & Suivi détaillé des licences, renouvellements et conformité & Perspective \\
\hline
Ticketing & Gestion des demandes support et incidents utilisateurs & Perspective \\
\hline
\captionsetup{justification=centering,margin=1cm}
\caption{Périmètre fonctionnel de la solution}
\end{longtable}
\section{État de l'art}

Avant de proposer notre solution, il est essentiel de réaliser une étude approfondie de l'existant sur le marché des outils de gestion de parc informatique. Cette analyse nous permettra de comprendre les forces et les faiblesses des solutions disponibles et de justifier notre choix de développer une application sur mesure.

\subsection{Étude des solutions existantes}

Plusieurs solutions de gestion de parc informatique (ITSM/CMDB) existent sur le marché, qu'elles soient open source ou commerciales. Nous analysons ci-dessous les principales solutions qui auraient pu répondre aux besoins d'ASTEELFLASH.

\subsubsection{GLPI (Gestionnaire Libre de Parc Informatique)}

GLPI est une solution open source française de gestion de parc informatique et de helpdesk, développée sous licence GPL. Elle permet la gestion complète de l'inventaire matériel et logiciel, la gestion des tickets d'incidents, le suivi des licences, la gestion financière des actifs et la planification des maintenances. GLPI bénéficie d'une large communauté d'utilisateurs et de développeurs, ce qui assure une bonne évolutivité et un support communautaire actif.

Cependant, GLPI présente certaines limites. Son interface utilisateur, bien que fonctionnelle, est parfois jugée peu intuitive et vieillissante. La personnalisation avancée nécessite des compétences techniques approfondies en PHP et MySQL. De plus, l'intégration native avec Microsoft Active Directory, bien que possible, requiert une configuration complexe et des plugins supplémentaires comme FusionInventory pour l'inventaire automatisé.

\subsubsection{OCS Inventory NG}

OCS Inventory NG (Open Computers and Software Inventory Next Generation) est un outil open source spécialisé dans l'inventaire automatique des ressources informatiques. Il fonctionne grâce à des agents installés sur les postes clients qui remontent automatiquement les informations matérielles et logicielles vers un serveur central. OCS Inventory se distingue par sa capacité de déploiement de paquets logiciels à distance et son architecture client-serveur performante.

Néanmoins, OCS Inventory est principalement un outil d'inventaire et ne propose pas nativement de fonctionnalités de gestion des affectations, des missions d'équipements ou de gestion des bons d'achat. Pour obtenir ces fonctionnalités, il est souvent couplé avec GLPI, ce qui complexifie l'architecture globale et la maintenance du système.

\subsubsection{Snipe-IT}

Snipe-IT est une application web open source dédiée à la gestion des actifs informatiques, développée en PHP avec le framework Laravel. Elle offre une interface moderne et responsive, un suivi complet des actifs avec attribution aux utilisateurs, une gestion des licences, des consommables et des accessoires, ainsi qu'un système d'audit trail (historique des modifications). Snipe-IT se distingue par sa facilité d'installation et son interface utilisateur moderne et intuitive.

Toutefois, Snipe-IT ne propose pas de fonctionnalités avancées telles que la gestion des missions temporaires d'équipements avec rappels automatiques par e-mail, ni de système d'alerte de stock minimum personnalisable. L'intégration avec Active Directory est disponible mais limitée en termes de gestion fine des groupes et des autorisations.

\subsubsection{Lansweeper}

Lansweeper est une solution commerciale de gestion des actifs IT qui offre une découverte automatique et sans agent des équipements réseau. Elle fournit un inventaire exhaustif du matériel, des logiciels, des licences et des utilisateurs. Lansweeper propose également des rapports prédéfinis et personnalisables, ainsi qu'une intégration native avec de nombreux systèmes tiers.

Cependant, Lansweeper est une solution payante dont le coût de licence peut être significatif pour une entreprise de taille moyenne. De plus, étant une solution propriétaire, la personnalisation est limitée aux options proposées par l'éditeur, ce qui ne permet pas de répondre à des besoins très spécifiques comme ceux d'ASTEELFLASH.

\subsection{Tableau comparatif des solutions}

Le tableau 1.2 présente une comparaison synthétique des solutions étudiées selon les critères pertinents pour ASTEELFLASH.

\begin{longtable}{|m{2.5cm}|m{2cm}|m{2cm}|m{2cm}|m{2cm}|m{2cm}|}
\hline
\textbf{Critère} & \textbf{GLPI} & \textbf{OCS Inv.} & \textbf{Snipe-IT} & \textbf{Lansweeper} & \textbf{Notre Solution} \\
\hline
\endhead
\endfoot
\endlastfoot
\hline
Open Source & Oui & Oui & Oui & Non & Oui \\
\hline
Intégration AD/LDAP & Partielle & Non & Partielle & Oui & Complète \\
\hline
Gestion Affectations & Basique & Non & Oui & Oui & Avancée \\
\hline
Gestion Missions & Non & Non & Non & Non & Oui \\
\hline
Alerte Stock Min. & Non & Non & Non & Non & Oui \\
\hline
E-mails automatiques & Plugin & Non & Limité & Oui & Oui \\
\hline
Tableau de Bord & Plugin & Non & Basique & Avancé & Personnalisé \\
\hline
Historique complet & Oui & Non & Oui & Oui & Oui \\
\hline
Personnalisation & Moyenne & Faible & Moyenne & Faible & Totale \\
\hline
Coût & Gratuit & Gratuit & Gratuit & Payant & Gratuit \\
\hline
\captionsetup{justification=centering,margin=1cm}
\caption{Tableau comparatif des solutions de gestion de parc informatique}
\end{longtable}

\subsection{Justification du développement sur mesure}

L'analyse comparative des solutions existantes met en évidence plusieurs constats importants qui justifient notre choix de développer une application sur mesure pour ASTEELFLASH :

\begin{itemize}[label=\textbullet,font=\normalsize]
    \item \textbf{Besoins spécifiques non couverts :} Aucune des solutions étudiées ne propose nativement la gestion des missions temporaires d'équipements avec un système de rappels automatiques par e-mail, fonctionnalité essentielle pour le fonctionnement d'ASTEELFLASH.
    \item \textbf{Intégration Active Directory complète :} Les solutions existantes offrent une intégration partielle ou limitée avec Active Directory. ASTEELFLASH nécessite une authentification LDAP complète avec une gestion fine des groupes d'utilisateurs et des autorisations basée sur l'annuaire d'entreprise.
    \item \textbf{Système d'alerte de stock minimum :} La fonctionnalité d'envoi automatique d'e-mails d'alerte lorsque le stock d'un type d'équipement atteint un seuil critique n'est pas disponible dans les solutions étudiées.
    \item \textbf{Personnalisation totale :} En développant une solution sur mesure, ASTEELFLASH conserve la maîtrise totale de l'application, de son évolution et de son adaptation aux processus internes de l'entreprise, sans dépendance vis-à-vis d'un éditeur tiers.
    \item \textbf{Maîtrise des coûts à long terme :} Bien que le développement initial représente un investissement, la solution sur mesure ne génère pas de coûts de licences récurrents et reste entièrement adaptable sans surcoûts liés à des modules complémentaires.
\end{itemize}

\section{ITIL et la gestion de parc informatique}

Notre projet s'inscrit dans le cadre des bonnes pratiques ITIL (Information Technology Infrastructure Library) pour la gestion des services informatiques.

\subsection{Présentation d'ITIL}

ITIL est un ensemble de bonnes pratiques pour la gestion des services informatiques (ITSM -- IT Service Management). Créé dans les années 1980 par l'agence centrale de l'informatique et des télécommunications du gouvernement britannique, ITIL est devenu la référence mondiale en matière de gestion des services IT. ITIL définit un cadre structuré pour planifier, délivrer et supporter les services informatiques au sein d'une organisation.

Les processus ITIL sont organisés autour de cinq phases du cycle de vie des services :
\begin{itemize}[label=\textbullet,font=\normalsize]
    \item \textbf{Stratégie des services :} Définir les services à offrir et les marchés à servir.
    \item \textbf{Conception des services :} Concevoir les services, les processus et les politiques nécessaires.
    \item \textbf{Transition des services :} Gérer les changements et les déploiements de services.
    \item \textbf{Exploitation des services :} Assurer la fourniture et le support quotidien des services.
    \item \textbf{Amélioration continue des services :} Mesurer et améliorer l'efficacité des services.
\end{itemize}

\subsection{Le concept de CMDB}

La CMDB (Configuration Management Database) est un composant central de la gestion des configurations selon ITIL. Il s'agit d'une base de données qui centralise l'ensemble des informations relatives aux éléments de configuration (CI -- Configuration Items) d'un système d'information. Un élément de configuration peut être un serveur, un poste de travail, un logiciel, un switch réseau, une imprimante ou tout autre composant du parc informatique.

La CMDB permet de :
\begin{itemize}[label=\textbullet,font=\normalsize]
    \item Recenser exhaustivement tous les actifs informatiques de l'entreprise.
    \item Tracer les relations et les dépendances entre les différents éléments de configuration.
    \item Fournir une base fiable pour la prise de décisions concernant les changements, les incidents et les problèmes.
    \item Supporter la gestion financière des actifs IT (coûts d'acquisition, amortissement, valeur résiduelle).
\end{itemize}

Notre application de gestion du parc informatique pour ASTEELFLASH s'inspire directement de ces concepts en implémentant les fonctionnalités essentielles d'une CMDB : inventaire centralisé, traçabilité des affectations, historique des changements et reporting statistique.

\section{Synthèse pour le choix de la méthodologie}

Dans cette section, nous allons présenter la méthodologie adoptée pour ce projet qui est la méthodologie agile qui conduit notre projet et garantit le bon déroulement des différentes phases.

La méthodologie agile est :
\begin{itemize}
   
\item [$-$] Une approche itérative et collaborative, capable de prendre en compte les besoins initiaux du client et ceux liés aux évolutions.
 \item [$-$]Nommée Manifeste Agile repose sur quatre grands principes :
\begin{itemize}[label=\textbullet]
    \item Priorité aux personnes et aux interactions sur les procédures et les outils.
    \item Priorité aux applications fonctionnelles sur une documentation abondante.
    \item Priorité de la collaboration avec le client sur la négociation de contrat
    \item Priorité de l’acceptation du changement sur la planification.
    \end{itemize}

    \begin{figure}[h!]
    \centering
    \includegraphics[scale=0.7]{Scrum.PNG}\\[0.5cm]
    \end{figure}\
    
\item [$-$]Basé sur un cycle de développement qui porte le client au centre, le client est impliqué dans la réalisation du début à la fin du projet.
L’implication du client dans le processus permet à l’équipe d’obtenir un feedback régulier afin d’appliquer directement les changements nécessaires. 
\item [$-$] Vise à accélérer le développement d’un logiciel. De plus, elle assure la réalisation d’un logiciel fonctionnel tout au long de la durée de sa création. 
En ce qui concerne la réalisation d’un logiciel, les instructions à répéter sont les suivantes :
\begin{itemize}[label=\textbullet]
    \item Les tests unitaires à la fin de chaque itération.
    \item Le développement de l’application web.
    \item L’intégration.
    \item La relecture et amélioration des codes [1].
    \end{itemize}
 
\end{itemize}
    \section{Scrum}

    La méthodologie SCRUM est utilisée par de nombreuses entreprises. De même, elle peut être utilisée pour développer des applications Web et mobile. Elle permet entre autres de développer des produits complexes et de produire la plus grande valeur métier dans la durée la plus courte. Elle permet aussi à ses utilisateurs d’être soudés pour accroitre la productivité. Pour conclure, cette méthodologie est la plus efficace qui nous permettra de réaliser notre projet dans un temps bien précis.
    \subsection{Principe du Scrum }
    Scrum en quelques mots :
    \begin{itemize}
    \item Scrum est un processus agile qui permet de produire la plus grande valeur métier dans la durée la plus courte.
    \item Du logiciel qui fonctionne est produit à chaque sprint (toutes les 2 à 4 semaines).
    \item Le métier définit les priorités, l’équipe s’organise elle-même pour déterminer la meilleure façon de produire les exigences les plus prioritaires.
    \item A chaque fin de sprint, tout le monde peut voir fonctionner le produit courant et décider soit de le livrer dans l’état, soit de continuer à l’améliorer pendant un sprint supplémentaire. [2]
    \end{itemize}

    \begin{figure}[H]
\centering
\frame{\includegraphics[width=1\columnwidth]{Scrum-process.png}}
\caption{Cycle de vie de la méthodologie SCRUM}
\label{fig:Scrum-process}
\end{figure}

    \subsection{Les rôles Scrum}

    D’après la méthode Scrum, chaque collaborateur a un rôle clairement défini dans la gestion du projet, elle définit seulement 3 rôles :
    \begin{itemize}
       
  
    \item [$-$] \textbf{Product Owner : }
    \begin{itemize}[label=\textbullet]
    \item Définit les fonctionnalités du produit.
    \item Choisit la date et le contenu de la release.
    \item Responsable du retour sur investissement.
    \item Définit les priorités dans le backlog en fonction de la valeur « métier ».
    \item Ajuste les fonctionnalités et les priorités à chaque sprint si nécessaire.
    \item Accepte ou rejette les résultats.
    \end{itemize}

    \item [$-$]\textbf{Scrum Master : }
     \begin{itemize}[label=\textbullet]
    \item Représente le management du projet.
    \item Responsable de faire appliquer par l’équipe les valeurs et les pratiques de Scrum.
    \item Élimine les obstacles.
    \item S’assure que l’équipe est complètement fonctionnelle et productive.
    \item Facilite une coopération poussée entre tous les rôles et fonctions.
    \item Protège l’équipe des interférences extérieures.
    \end{itemize}

  \item [$-$] \textbf{Equipe de Développement: }
     \begin{itemize}[label=\textbullet,font=\normalsize]
    \item Chargée de transformer les besoins exprimés par le Product Owner en fonctionnalités utilisables. Elle est pluridisciplinaire et peut donc encapsuler d’autres rôles tels que développeur, architecte logiciel, DBA, analyste fonctionnel, graphiste/ergonome, ingénieur système.
    \item L’équipe s’organise par elle-même.
    \item La composition de l’équipe ne doit pas changer pendant un Sprint [3].
    \end{itemize}
  \end{itemize}
  \subsection{ Backlog de produit :}
     Le backlog est élaboré avant le lancement des sprints, dans la phase de préparation (ou sprint 0). Il est utilisé pour planifier la release, puis à chaque sprint, lors de la réunion de planification du sprint pour décider du sous-ensemble qui sera réalisé. C’est donc un outil essentiel pour la planification. Mais il est aussi, par sa nature, un maillon de la gestion des exigences, puisqu’on y collecte ce que doit faire le produit. [4]


  \subsection{Adaptation de Scrum au contexte du projet}

Dans notre projet, Scrum a été adapté au contexte d'un stage de fin d'études réalisé au sein d'une entreprise industrielle. L'équipe étant réduite, certaines cérémonies ont été allégées tout en conservant l'esprit agile : échanges fréquents avec le Product Owner, priorisation continue du backlog, démonstration des incréments et intégration progressive des remarques.

La démarche suivie s'articule autour des points suivants :

\begin{itemize}[label=\textbullet,font=\normalsize]
    \item \textbf{Réunions de cadrage :} elles ont permis de comprendre le fonctionnement actuel du département IT, les irritants rencontrés et les fonctionnalités attendues.
    \item \textbf{Planification des sprints :} chaque sprint a été associé à un objectif clair afin d'éviter la dispersion et de produire un résultat vérifiable.
    \item \textbf{Validation régulière :} les interfaces développées ont été présentées progressivement pour recueillir les remarques du Product Owner et ajuster les priorités.
    \item \textbf{Gestion des changements :} les demandes apparues en cours de projet ont été intégrées lorsque leur valeur métier était importante et lorsque leur impact sur le planning restait maîtrisé.
    \item \textbf{Amélioration continue :} les anomalies détectées durant les tests ont été corrigées au fil de l'eau afin d'améliorer la qualité de l'application avant la livraison finale.
\end{itemize}

Cette adaptation a permis de conserver une communication directe avec l'entreprise, d'éviter l'effet tunnel et de livrer progressivement les modules essentiels : authentification, gestion des équipements, suivi des affectations, missions, commandes et fournisseurs.
\section{Conclusion}
Dans ce premier chapitre, nous avons commencé par présenter l'organisme d'accueil ASTEELFLASH, en détaillant son historique, ses activités, ses certifications et son organigramme. Nous avons ensuite réalisé une étude de l'existant qui a mis en évidence les limites de la gestion actuelle du parc informatique basée sur des fichiers Excel et des décharges papier.

L'analyse de la problématique nous a conduit à proposer le développement d'une application web dédiée à la gestion du parc informatique et CMDB. Une étude comparative approfondie des solutions existantes sur le marché (GLPI, OCS Inventory, Snipe-IT, Lansweeper) a démontré qu'aucune d'entre elles ne répond pleinement aux besoins spécifiques d'ASTEELFLASH, notamment en termes d'intégration avec Active Directory, de gestion des missions temporaires et d'alertes de stock minimum.

Nous avons également situé notre projet dans le cadre des bonnes pratiques ITIL et du concept de CMDB, avant de présenter la méthodologie agile Scrum adoptée pour le développement. Cette méthodologie nous permettra de livrer des incréments fonctionnels à chaque sprint, en impliquant le client tout au long du processus.

Dans le chapitre suivant, nous allons entamer la spécification des besoins fonctionnels et non fonctionnels, présenter l'équipe Scrum, le backlog du produit et les diagrammes de cas d'utilisation et de classes.

